% Figure 6-1 -- The context build pipeline.
% Restyled from context-pipeline.gv; the ten stages, their order and numbers,
% the three margin annotations and the manifest note shape are unchanged.
%
% Original element -> restyled element
%   s1  "1 . Parse request"          -> slate card; the number "1" is a badge, the words are the card text
%   s2  "2 . Candidate sources"      -> cyan card, badge 2
%   s3  "3 . Authorize"              -> red card with a heavier border (as in the source), badge 3
%   s4  "4 . Fetch pinned versions"  -> cyan card, badge 4
%   s5  "5 . Classify and sanitize"  -> cyan card, badge 5
%   s6  "6 . Deduplicate"            -> cyan card, badge 6
%   s7  "7 . Rank"                   -> cyan card, badge 7
%   s8  "8 . Allocate tokens"        -> cyan card, badge 8
%   s9  "9 . Ordered references"     -> blue card, badge 9
%   s10 "10 . ContextManifest"       -> amber card with a folded corner (the source's note shape), badge 10
%   mdl "model request"              -> teal card
%   n3  "before a single byte is fetched"          -> red annotation panel beside stage 3
%   n9  "references only - no content yet"         -> slate annotation panel beside stage 9
%   nm  "the only boundary that materializes bytes" -> teal annotation panel beside "model request"
%   edges: s1>s2>s3>s4>s5>s6>s7>s8>s9>s10>mdl (10 solid arrows, same direction)
\documentclass[border=6pt]{standalone}
\usepackage{professional-diagrams}

% ---- local constructs (candidates for promotion to professional-diagrams.sty) ----
% num=<n>: a round number badge at the left of a card (same slot as icon=...).
% \dgdoccorner{node}{role}: the folded top-right corner of a note-shaped card.
\tikzset{
  num/.style={execute at begin node={\leftskip=2.75em\relax},
    append after command={
      node[circle, fill=dg-\dgcurrentrole-line, text=white, minimum size=19pt,
           inner sep=0pt, font=\dglabelsize]
        at ($(\tikzlastnode.west)+(\dgpadx+0.95em,0)$) {#1}}},
  note tag/.style={callout=#1, minimum height=33pt, text=dg-#1-ink, text width=170pt, font=\dgdescsize},
  leader/.style={draw=dg-#1-line!70, line width=0.8pt, line cap=round},
}

\newcommand{\dgdoccorner}[2]{%
  \path[fill=white] ([xshift=-11pt]#1.north east) -- ([xshift=1pt,yshift=1pt]#1.north east)
    -- ([yshift=-11pt]#1.north east) -- cycle;
  \path[fill=dg-#2-pill, draw=dg-#2-line, line width=\dgstroke, line join=round]
    ([xshift=-11pt]#1.north east) -- ([xshift=-11pt,yshift=-11pt]#1.north east)
    -- ([yshift=-11pt]#1.north east) -- cycle;}

\begin{document}
\begin{tikzpicture}
  % ---- stages (centre x = 130, pitch 49) ----
  \node[dg card=slate, num=1, text width=225pt, minimum height=33pt] (s1) at (130,0)
    {\dgcontent{Parse request}{}{}};
  \node[dg card=cyan, num=2, text width=225pt, minimum height=33pt] (s2) at (130,-49)
    {\dgcontent{Candidate sources}{}{}};
  \node[dg card=red, num=3, text width=225pt, minimum height=33pt, line width=1.8pt] (s3) at (130,-98)
    {\dgcontent{Authorize}{}{}};
  \node[dg card=cyan, num=4, text width=225pt, minimum height=33pt] (s4) at (130,-147)
    {\dgcontent{Fetch pinned versions}{}{}};
  \node[dg card=cyan, num=5, text width=225pt, minimum height=33pt] (s5) at (130,-196)
    {\dgcontent{Classify and sanitize}{}{}};
  \node[dg card=cyan, num=6, text width=225pt, minimum height=33pt] (s6) at (130,-245)
    {\dgcontent{Deduplicate}{}{}};
  \node[dg card=cyan, num=7, text width=225pt, minimum height=33pt] (s7) at (130,-294)
    {\dgcontent{Rank}{}{}};
  \node[dg card=cyan, num=8, text width=225pt, minimum height=33pt] (s8) at (130,-343)
    {\dgcontent{Allocate tokens}{}{}};
  \node[dg card=blue, num=9, text width=225pt, minimum height=33pt] (s9) at (130,-392)
    {\dgcontent{Ordered references}{}{}};
  \node[dg card=amber, num=10, text width=225pt, minimum height=33pt] (s10) at (130,-441)
    {\dgcontent{ContextManifest}{}{}};
  \node[dg card=teal, num={\dgicon{paper-plane}}, text width=225pt, minimum height=33pt] (mdl) at (130,-490)
    {\dgcontent{Model request}{}{}};

  \dgdoccorner{s10}{amber}

  % ---- annotations beside the three load-bearing steps ----
  \node[note tag=red, anchor=west] (n3) at (300,-98)
    {Before a single byte is fetched};
  \node[note tag=slate, anchor=west] (n9) at (300,-392)
    {References only \textemdash\\no content yet};
  \node[note tag=teal, anchor=west] (nm) at (300,-490)
    {The only boundary that materializes bytes};
  \draw[leader=red] (s3.east) -- (n3.west);
  \draw[leader=slate] (s9.east) -- (n9.west);
  \draw[leader=teal] (mdl.east) -- (nm.west);

  % ---- flow ----
  \foreach \a/\b in {s1/s2, s2/s3, s3/s4, s4/s5, s5/s6, s6/s7, s7/s8, s8/s9, s9/s10, s10/mdl}
    \draw[flow in=blue] (\a.south) -- (\b.north);
\end{tikzpicture}
\end{document}
